\documentclass[11pt,a4paper]{article}
\usepackage{shared/parscover}

% ============================================================================
% Dual-Layer Architecture: EVM Rollup with Privacy-Preserving Session Layer
% Pars Network Technical Whitepaper PNP-001
% ============================================================================

\usepackage[utf8]{inputenc}
\usepackage[T1]{fontenc}
\usepackage{amsmath,amssymb,amsthm}
\usepackage{graphicx}
\usepackage{hyperref}
\usepackage{algorithm}
\usepackage{algpseudocode}
\usepackage{booktabs}
\usepackage{listings}
\input{shared/lstlang}
\usepackage{xcolor}
\usepackage{tikz}
\usepackage{enumitem}
\usepackage{caption}
\usepackage{subcaption}
\usepackage{fancyhdr}
\usepackage{abstract}

\geometry{margin=1in}

\hypersetup{
    colorlinks=true,
    linkcolor=blue!70!black,
    citecolor=green!50!black,
    urlcolor=blue!60!black
}

\lstset{
    basicstyle=\ttfamily\small,
    breaklines=true,
    frame=single,
    backgroundcolor=\color{gray!10},
    keywordstyle=\color{blue!70!black}\bfseries,
    commentstyle=\color{green!50!black}\itshape,
    stringstyle=\color{red!60!black},
    numbers=left,
    numberstyle=\tiny\color{gray},
    numbersep=5pt
}

\usetikzlibrary{arrows.meta,positioning,shapes.geometric,fit,calc}

\newtheorem{definition}{Definition}[section]
\newtheorem{theorem}{Theorem}[section]
\newtheorem{lemma}{Lemma}[section]
\newtheorem{proposition}{Proposition}[section]
\newtheorem{corollary}{Corollary}[section]

\pagestyle{fancy}
\fancyhf{}
\fancyhead[L]{\small Cyrus Pahlavi, Pars Team}
\fancyhead[R]{\small Dual-Layer Architecture}
\fancyfoot[C]{\thepage}

% ============================================================================
\title{%
    \textbf{Dual-Layer Architecture: EVM Rollup with\\
    Privacy-Preserving Session Layer}\\[0.5em]
    \large Pars Network Technical Whitepaper PNP-001
}

\author{%
    Cyrus Pahlavi, Pars Team\\
    \texttt{research@pars.network}\\[0.5em]
    \small Version 1.0 --- February 2026
}

\date{}

% ============================================================================
\begin{document}
\parscoverpage

\thispagestyle{fancy}

% ============================================================================
\begin{abstract}
\noindent
We present the Pars dual-layer architecture, a system that couples an
EVM-compatible optimistic rollup (Layer~2) with a privacy-preserving session
layer (Layer~S) to deliver sovereign infrastructure for the Persian diaspora.
The L2 rollup inherits Ethereum's security guarantees while providing
low-cost, high-throughput transaction processing.  The session layer
operates as an overlay network of daemon processes that manage end-to-end
encrypted communication channels, ephemeral state, and privacy-preserving
metadata routing.  Cross-layer communication is achieved through a set of
bridge contracts and a \emph{session anchor protocol} that commits session
state roots to the rollup without revealing session contents.  We formalise
the state synchronisation protocol, prove its consistency under asynchronous
network conditions, and present benchmarks demonstrating sub-second finality
for session operations with amortised gas costs below 500~gas units per
session commitment.  The architecture enables censorship-resistant
communication, private governance participation, and financial sovereignty
for communities operating under authoritarian surveillance regimes.
\end{abstract}

\vspace{1em}
\noindent\textbf{Keywords:} rollup architecture, EVM compatibility, privacy
layer, session management, diaspora infrastructure, censorship resistance

\tableofcontents
\newpage

% ============================================================================
\section{Introduction}
\label{sec:introduction}

The Persian diaspora spans over eight million individuals across more than
forty countries~\cite{amanat2012persian}.  Members require digital services
that are simultaneously accessible from jurisdictions with varying degrees
of internet freedom, resistant to state-level censorship and surveillance,
and capable of preserving cultural and economic connections across borders.

Existing blockchain platforms offer either strong financial guarantees with
transparent ledgers (Ethereum, Bitcoin) or privacy-focused transactions with
limited programmability (Zcash, Monero).  Neither paradigm alone addresses
the compound requirements of a diaspora community: the need for programmable
financial instruments \emph{and} private communication \emph{and} censorship
resistance \emph{and} low operational costs.

\subsection{Problem Statement}

We identify four fundamental requirements for diaspora infrastructure:

\begin{enumerate}[label=(\roman*)]
    \item \textbf{Financial Sovereignty:} The ability to hold, transfer, and
    programme digital assets without dependence on any single jurisdiction's
    banking infrastructure.

    \item \textbf{Communication Privacy:} End-to-end encrypted messaging and
    data exchange with forward secrecy, resistant to traffic analysis and
    metadata correlation.

    \item \textbf{Censorship Resistance:} Infrastructure that remains
    operational even when state actors attempt to block, throttle, or
    intercept network traffic.

    \item \textbf{Cost Efficiency:} Transaction and communication costs
    accessible to individuals in economically constrained environments.
\end{enumerate}

\subsection{Contribution}

This paper presents the Pars dual-layer architecture, which addresses all
four requirements through a coupling of two complementary layers:

\begin{itemize}
    \item \textbf{Layer~2 (L2):} An EVM-compatible optimistic rollup that
    provides programmable financial infrastructure with Ethereum-grade
    security at a fraction of the cost.

    \item \textbf{Layer~S (Session Layer):} A privacy-preserving overlay
    network that manages encrypted communication sessions, ephemeral state,
    and metadata-resistant routing.
\end{itemize}

The key innovation is the \emph{session anchor protocol}, which allows the
session layer to commit cryptographic proofs of session state to the rollup
without revealing session contents, enabling verifiable privacy.

\subsection{Paper Organisation}

Section~\ref{sec:background} reviews related work.
Section~\ref{sec:l2design} details the L2 rollup.
Section~\ref{sec:session} presents the session layer.
Section~\ref{sec:crosslayer} describes cross-layer communication.
Section~\ref{sec:statesync} formalises state synchronisation.
Section~\ref{sec:security} provides security analysis.
Section~\ref{sec:performance} presents benchmarks.
Section~\ref{sec:governance} discusses governance.
Section~\ref{sec:conclusion} concludes.

% ============================================================================
\section{Background and Related Work}
\label{sec:background}

\subsection{Optimistic Rollups}

Optimistic rollups execute transactions off-chain and post compressed data
to Ethereum L1, relying on fraud proofs for
security~\cite{kalodner2018arbitrum,optimism2021}.  Only disputed state
transitions require L1 verification, amortising costs across many
transactions.  Arbitrum~\cite{kalodner2018arbitrum} and
Optimism~\cite{optimism2021} achieve 10--100$\times$ cost reduction relative
to L1 execution.

\subsection{Privacy-Preserving Networks}

Tor~\cite{dingledine2004tor} provides anonymous communication through onion
routing but suffers from latency and vulnerability to traffic correlation by
global adversaries.  I2P~\cite{zantout2011i2p} offers garlic routing with
better resistance to timing attacks but lower adoption.
Mixnets~\cite{chaum1981untraceable,danezis2003mixminion} provide stronger
guarantees through message batching and reordering at the cost of
significant latency.

\subsection{Session Management Protocols}

The Signal Protocol~\cite{marlinspike2016signal} provides end-to-end
encryption with forward secrecy through the Double Ratchet algorithm.
Matrix~\cite{hodgson2019matrix} extends this to federated group
communication.  Briar~\cite{briar2017} adds mesh networking capability for
environments where internet access is unreliable or monitored.

\subsection{Privacy Layers for Blockchains}

Aztec~\cite{williamson2018aztec} implements a privacy layer for Ethereum
using zk-SNARKs.  Tornado Cash~\cite{pertsev2019tornado} provides
transaction mixing through zero-knowledge proofs.
Penumbra~\cite{penumbra2022} builds a privacy-focused chain within the
Cosmos ecosystem.  None integrate communication privacy with financial
privacy in a unified architecture.

\subsection{Positioning of Pars}

Pars occupies a unique position by integrating an L2 rollup for financial
operations with a session layer for communication, connected through
cryptographic anchoring.  Users transact privately, communicate securely,
and participate in governance within a single coherent infrastructure.

% ============================================================================
\section{Layer~2 Rollup Architecture}
\label{sec:l2design}

\subsection{Design Principles}

The Pars L2 is designed around four principles:

\begin{enumerate}
    \item \textbf{Full EVM Equivalence:} Byte-for-byte compatibility with
    the Ethereum Virtual Machine, enabling deployment of existing Solidity
    contracts without modification.

    \item \textbf{Optimistic Execution:} Transactions are presumed valid
    with a challenge period for fraud detection.

    \item \textbf{Data Availability on L1:} All transaction data is posted
    to Ethereum L1 as calldata or blob data via
    EIP-4844~\cite{eip4844}, ensuring that anyone can reconstruct the L2
    state.

    \item \textbf{Session Awareness:} The rollup includes native precompiles
    for session layer interaction.
\end{enumerate}

\subsection{Rollup Components}

\subsubsection{Sequencer}

The sequencer orders transactions and produces L2 blocks.  Pars operates a
decentralised sequencer set using a weighted round-robin selection:

\begin{definition}[Sequencer Set]
    Let $\mathcal{S} = \{s_1, \dots, s_n\}$ be the set of registered
    sequencers with stake weights $\{w_1, \dots, w_n\}$ where
    $\sum_{i=1}^n w_i = 1$.  The sequencer for block $b$ is:
    \[
        \mathrm{seq}(b) = s_j, \quad
        j = \min\Bigl\{k : \sum_{i=1}^k w_i
            > \frac{b \bmod P}{P}\Bigr\},
    \]
    where $P$ is the rotation period in blocks.
\end{definition}

\subsubsection{State Commitment}

The rollup maintains a Merkle-Patricia trie for account state.  State roots
are committed to L1 at regular intervals:

\begin{definition}[State Commitment]
    A state commitment $C_b$ for block $b$ consists of:
    \[
        C_b = \bigl\langle
            \mathrm{root}(\mathcal{T}_b),\;
            H(\mathrm{txdata}_b),\;
            \mathrm{root}(\mathcal{R}_b),\;
            \mathrm{root}(\mathcal{A}_b)
        \bigr\rangle,
    \]
    where $\mathcal{T}_b$ is the state trie, $\mathrm{txdata}_b$ is the
    compressed transaction batch, $\mathcal{R}_b$ is the receipt trie, and
    $\mathcal{A}_b$ is the session anchor trie.
\end{definition}

Inclusion of the session anchor trie $\mathcal{A}_b$ is the key extension
beyond standard rollup designs.

\subsubsection{Fraud Proofs}

Pars employs interactive fraud proofs following the Arbitrum
model~\cite{kalodner2018arbitrum}.  When a validator disputes a state
commitment, a bisection protocol narrows the disagreement to a single
instruction execution, verified on L1.

\begin{algorithm}[H]
\caption{Interactive Fraud Proof Protocol}
\label{alg:fraud-proof}
\begin{algorithmic}[1]
\Require State commitments $C_a, C_b$ where $a < b$ and $C_b$ is disputed
\State $\mathit{lo} \gets a$,\; $\mathit{hi} \gets b$
\While{$\mathit{hi} - \mathit{lo} > 1$}
    \State $\mathit{mid} \gets \lfloor (\mathit{lo} + \mathit{hi}) / 2 \rfloor$
    \If{Challenger agrees with $C_{\mathit{mid}}$}
        \State $\mathit{lo} \gets \mathit{mid}$
    \Else
        \State $\mathit{hi} \gets \mathit{mid}$
    \EndIf
\EndWhile
\State Verify single-step execution from $C_{\mathit{lo}}$ to
       $C_{\mathit{hi}}$ on L1
\State Slash the party that produced the incorrect commitment
\end{algorithmic}
\end{algorithm}

\subsection{Session-Aware Precompiles}

The Pars L2 extends the EVM with four precompiled contracts:

\begin{table}[h]
\centering
\caption{Session-Aware Precompile Addresses}
\label{tab:precompiles}
\begin{tabular}{@{}lll@{}}
\toprule
\textbf{Address} & \textbf{Name} & \textbf{Function} \\
\midrule
\texttt{0x0100} & \texttt{SESSION\_ANCHOR}  & Commit session state root \\
\texttt{0x0101} & \texttt{SESSION\_VERIFY}  & Verify session proof \\
\texttt{0x0102} & \texttt{SESSION\_RESOLVE} & Resolve session identity \\
\texttt{0x0103} & \texttt{SESSION\_BRIDGE}  & Bridge assets to session layer \\
\bottomrule
\end{tabular}
\end{table}

\subsubsection{Session Anchor Precompile}

\begin{lstlisting}[language=Solidity,caption={Session Anchor Interface}]
interface ISessionAnchor {
    /// @notice Commit a session state root with validity proof
    /// @param sessionRoot Merkle root of session state
    /// @param proof       ZK proof of valid session state transition
    /// @param epoch       Session epoch number
    function anchor(
        bytes32 sessionRoot,
        bytes calldata proof,
        uint256 epoch
    ) external returns (bool);

    /// @notice Verify session state inclusion
    /// @param sessionRoot Committed session root
    /// @param leaf        Leaf value to verify
    /// @param proof       Merkle inclusion proof
    function verify(
        bytes32 sessionRoot,
        bytes32 leaf,
        bytes32[] calldata proof
    ) external view returns (bool);
}
\end{lstlisting}

\subsection{Transaction Batching and Compression}

Transaction data is compressed before posting to L1 using domain-specific
compression:

\begin{enumerate}
    \item \textbf{Address Table:} Frequently used addresses are replaced
    with 2-byte indices, reducing 20-byte addresses to 2~bytes for common
    interactions.

    \item \textbf{Delta Encoding:} Nonce and value fields are delta-encoded
    relative to the previous transaction in the batch.

    \item \textbf{Signature Aggregation:} BLS signature
    aggregation~\cite{boneh2001bls} reduces $n$ signatures to a single
    aggregate, saving $64(n-1)$~bytes per batch.

    \item \textbf{Zlib Compression:} The encoded batch receives a final
    zlib pass at level~9.
\end{enumerate}

Empirically this achieves 8--12$\times$ compression, reducing per-transaction
L1 cost to approximately 200--400~gas.

\subsection{EIP-4844 Blob Integration}

With EIP-4844 blob transactions~\cite{eip4844}, Pars preferentially uses
blob space for data availability:

\begin{proposition}[Cost Reduction via Blobs]
    Let $c_{\mathrm{calldata}} = 16$~gas/byte and
    $c_{\mathrm{blob}} \approx 1$~gas/byte.
    For a batch of $B$~compressed bytes the cost ratio is
    $c_{\mathrm{blob}} / c_{\mathrm{calldata}} = 1/16$,
    yielding an additional $16\times$ cost reduction.
\end{proposition}

% ============================================================================
\section{Session Layer Design}
\label{sec:session}

\subsection{Overview}

The session layer (Layer~S) is an overlay network of daemon processes that
run alongside L2 nodes.

\begin{definition}[Session Layer]
    $\mathcal{L}_S = (\mathcal{N}, \mathcal{C}, \mathcal{R}, \mathcal{E})$
    where $\mathcal{N}$ is the set of session daemons, $\mathcal{C}$ the
    active channels, $\mathcal{R}$ the routing protocol, and
    $\mathcal{E}$ the epoch management system.
\end{definition}

\subsection{Session Daemon Architecture}

\subsubsection{Key Manager}

Each daemon maintains:

\begin{itemize}
    \item \textbf{Identity Key Pair:} Long-term Ed25519 for daemon
    authentication.
    \item \textbf{Session Keys:} Ephemeral X25519 pairs for
    Diffie-Hellman exchange.
    \item \textbf{Prekey Bundle:} One-time prekeys (OTPKs) uploaded to the
    session directory for asynchronous session establishment.
    \item \textbf{Ratchet State:} Double
    Ratchet~\cite{marlinspike2016signal} state for each active session.
\end{itemize}

\subsubsection{Channel Manager}

\begin{algorithm}[H]
\caption{Channel Establishment Protocol}
\label{alg:channel-establish}
\begin{algorithmic}[1]
\Require Initiator $A$, Responder $B$
\State $A$ retrieves $B$'s prekey bundle from directory
\State $A$ performs X3DH key agreement~\cite{marlinspike2016x3dh}:
    \Statex \quad $\mathrm{SK} \gets
        \mathrm{X3DH}(\mathrm{IK}_A, \mathrm{EK}_A,
        \mathrm{IK}_B, \mathrm{SPK}_B, \mathrm{OPK}_B)$
\State $A$ initialises Double Ratchet with SK
\State $A$ sends initial message with $\mathrm{EK}_A$ and used
       $\mathrm{OPK}_B$ identifier
\State $B$ receives, reconstructs SK via X3DH
\State $B$ initialises Double Ratchet with SK
\State Channel $c_{AB}$ established with forward secrecy
\end{algorithmic}
\end{algorithm}

\subsubsection{Routing Engine}

The routing engine implements a modified Sphinx~\cite{danezis2009sphinx}
packet format:

\begin{definition}[Sphinx Packet]
    $P = (\alpha, \beta, \gamma, \delta)$ where $\alpha$ is a group
    element for key derivation, $\beta$ the encrypted routing header,
    $\gamma$ a MAC, and $\delta$ the encrypted payload.
\end{definition}

Messages traverse a minimum of three relay nodes drawn from distinct
availability zones.

\subsection{Ephemeral State Management}

The session layer maintains ephemeral state never persisted to L2:

\begin{itemize}
    \item Active channel states and ratchet keys
    \item Pending message queues for offline recipients
    \item Relay routing tables
    \item Reputation scores for relay nodes
\end{itemize}

Communication state requires only eventual delivery and forward secrecy,
not global consistency.

\subsection{Session Epochs}

The session layer operates in epochs (default 1~hour).  At each boundary:

\begin{enumerate}
    \item All active sessions perform a ratchet step, refreshing keys.
    \item Each daemon computes a Merkle root of its local state.
    \item Epoch roots are aggregated into a global session state root.
    \item The global root is committed to L2 via the session anchor
    precompile.
\end{enumerate}

\begin{theorem}[Forward Secrecy Guarantee]
    \label{thm:forward-secrecy}
    If a daemon's long-term key is compromised at time $t$, messages sent
    before epoch $e < \lfloor t / T_{\mathrm{epoch}} \rfloor$ remain
    confidential, provided at least one intermediate ratchet step used an
    uncompromised ephemeral key.
\end{theorem}

\begin{proof}
    Each ratchet step derives a new chain key
    $\mathrm{CK}_i = \mathrm{KDF}(\mathrm{CK}_{i-1},
    \mathrm{DH}(ek_i^A, ek_i^B))$.
    Compromising $\mathrm{IK}_A$ does not reveal any $ek_i^A$ for
    $i < \lfloor t / T_{\mathrm{epoch}} \rfloor$ because ephemeral keys
    are deleted after use.  Therefore $\mathrm{CK}_i$ cannot be
    reconstructed and messages encrypted under derived keys remain
    confidential.
\end{proof}

% ============================================================================
\section{Cross-Layer Communication}
\label{sec:crosslayer}

\subsection{Communication Model}

Cross-layer communication follows an asymmetric model:

\begin{itemize}
    \item \textbf{L2 $\to$ Session:} The rollup emits events that session
    daemons monitor (publish-subscribe).

    \item \textbf{Session $\to$ L2:} The session layer commits state roots
    and proofs via the anchor precompile (batched L2 transaction).
\end{itemize}

\subsection{Session Anchor Protocol}

\begin{definition}[Session Anchor]
    A session anchor $A_e$ for epoch $e$ is
    $\langle e, r_e, \pi_e, \sigma_e \rangle$
    where $r_e$ is the session state Merkle root, $\pi_e$ is a
    zero-knowledge proof of valid transition from $r_{e-1}$ to $r_e$, and
    $\sigma_e$ is an aggregate BLS signature from the session committee.
\end{definition}

\begin{algorithm}[H]
\caption{Session Anchor Commitment}
\label{alg:anchor-commit}
\begin{algorithmic}[1]
\Require Session committee $\mathcal{S}_c$, previous anchor $A_{e-1}$
\State Each $n_i \in \mathcal{S}_c$ computes local root $r_i$
\State Coordinator aggregates:
       $r_e \gets \mathrm{MerkleRoot}(\{r_i\})$
\State Generate ZK proof:
       $\pi_e \gets \mathrm{Prove}(r_{e-1}, r_e, \mathrm{transitions})$
\State Each node signs:
       $\sigma_i \gets \mathrm{Sign}(\mathrm{sk}_i, (e, r_e, \pi_e))$
\State Aggregate:
       $\sigma_e \gets \mathrm{BLSAggregate}(\{\sigma_i\})$
\State Submit $A_e$ to \texttt{SESSION\_ANCHOR} precompile
\State L2 verifies $\pi_e$ and $\sigma_e$; stores $A_e$ in anchor trie
\end{algorithmic}
\end{algorithm}

\subsection{Bridge Contracts}

\begin{lstlisting}[language=Solidity,caption={Session Bridge Contract}]
contract SessionBridge {
    mapping(bytes32 => uint256) public locked;
    mapping(bytes32 => bool) public processed;

    event Deposit(
        bytes32 indexed sessionId,
        address indexed depositor,
        uint256 amount
    );
    event Withdrawal(
        bytes32 indexed sessionId,
        address indexed recipient,
        uint256 amount,
        bytes32 proofRoot
    );

    function deposit(bytes32 sid) external payable {
        locked[sid] += msg.value;
        emit Deposit(sid, msg.sender, msg.value);
    }

    function withdraw(
        bytes32 sid,
        address to,
        uint256 amt,
        bytes32[] calldata proof,
        bytes32 anchorRoot
    ) external {
        bytes32 leaf = keccak256(
            abi.encodePacked(sid, to, amt)
        );
        require(
            _verifyAnchor(anchorRoot, leaf, proof),
            "Bad proof"
        );
        bytes32 wid = keccak256(
            abi.encodePacked(sid, amt, block.number)
        );
        require(!processed[wid], "Replayed");
        processed[wid] = true;
        locked[sid] -= amt;
        payable(to).transfer(amt);
        emit Withdrawal(sid, to, amt, anchorRoot);
    }

    function _verifyAnchor(
        bytes32 root,
        bytes32 leaf,
        bytes32[] calldata proof
    ) internal view returns (bool) {
        (bool ok, bytes memory res) =
            address(0x0101).staticcall(
                abi.encode(root, leaf, proof)
            );
        return ok && abi.decode(res, (bool));
    }
}
\end{lstlisting}

\subsection{Event Subscription}

\begin{lstlisting}[language=Go,caption={Session Event Subscriber}]
type EventSubscriber struct {
    l2       *ethclient.Client
    daemon   *SessionDaemon
    topics   []common.Hash
}

func (s *EventSubscriber) Run(
    ctx context.Context,
) error {
    q := ethereum.FilterQuery{
        Addresses: []common.Address{bridgeAddr},
        Topics:    [][]common.Hash{s.topics},
    }
    ch := make(chan types.Log, 64)
    sub, err := s.l2.SubscribeFilterLogs(ctx, q, ch)
    if err != nil {
        return fmt.Errorf("subscribe: %w", err)
    }
    defer sub.Unsubscribe()
    for {
        select {
        case log := <-ch:
            s.daemon.HandleL2Event(log)
        case err := <-sub.Err():
            return fmt.Errorf("sub err: %w", err)
        case <-ctx.Done():
            return nil
        }
    }
}
\end{lstlisting}

% ============================================================================
\section{State Synchronisation}
\label{sec:statesync}

\subsection{Consistency Model}

\begin{itemize}
    \item \textbf{L2 State:} Strong consistency via total ordering by the
    sequencer.  Sub-second soft finality; 7-day L1 finality.

    \item \textbf{Session State:} Eventual consistency within epochs, with
    causal ordering within individual channels.
\end{itemize}

\begin{definition}[Causal Consistency]
    For events $e_1, e_2$ in the session layer, if $e_1 \to e_2$ then
    every correct node observing $e_2$ has previously observed $e_1$.
\end{definition}

\subsection{Synchronisation Protocol}

Three phases per epoch:

\begin{enumerate}
    \item \textbf{Collection:} Daemons accumulate state updates and
    maintain local Merkle trees.
    \item \textbf{Aggregation:} The session committee aggregates local
    roots via a binary-tree protocol.
    \item \textbf{Commitment:} The global root is committed to L2 via
    the session anchor protocol.
\end{enumerate}

\begin{theorem}[Synchronisation Consistency]
    Under $f < n/3$ Byzantine session committee members, the protocol
    produces a consistent session anchor within $O(\log n)$ rounds.
\end{theorem}

\begin{proof}
    The aggregation uses a binary tree with $\lceil \log_2 n \rceil$
    levels.  At each level, pairs exchange and verify Merkle roots.  With
    $f < n/3$ Byzantine nodes, at least $2n/3$ nodes at each level produce
    correct intermediate roots.  Any inconsistency is detected by the
    sibling through root verification.  After $\lceil \log_2 n \rceil$
    rounds the root node holds a consistent global root, signed by at
    least $2n/3 + 1$ committee members whose aggregate BLS signature is a
    compact proof of agreement.
\end{proof}

\subsection{Conflict Resolution}

\begin{definition}[Session Vector Clock]
    Each daemon $d_i$ maintains $V_i = [v_1, \dots, v_n]$ where $v_j$
    represents $d_i$'s knowledge of $d_j$'s latest state version.
\end{definition}

For concurrent updates (neither $V_a \leq V_b$ nor $V_b \leq V_a$), the
update with the lexicographically smaller session ID takes precedence.

\subsection{Checkpoint and Recovery}

\begin{algorithm}[H]
\caption{Checkpoint and Recovery}
\label{alg:checkpoint}
\begin{algorithmic}[1]
\Procedure{Checkpoint}{daemon $d$, epoch $e$}
    \State $\mathit{state} \gets d.\mathrm{Serialise}()$
    \State $k \gets \mathrm{KDF}(d.\mathrm{master}, e)$
    \State $\mathit{ct} \gets \mathrm{AES\text{-}GCM}(k, \mathit{state})$
    \State $\mathit{mac} \gets \mathrm{HMAC}(k, \mathit{ct})$
    \State Write $(\mathit{ct}, \mathit{mac})$ to local storage
\EndProcedure
\Procedure{Recover}{daemon $d$, epoch $e$}
    \State Read $(\mathit{ct}, \mathit{mac})$ from storage
    \State $k \gets \mathrm{KDF}(d.\mathrm{master}, e)$
    \State Verify $\mathrm{HMAC}(k, \mathit{ct}) = \mathit{mac}$
    \State $\mathit{state} \gets \mathrm{AES\text{-}GCM\text{-}Dec}(k, \mathit{ct})$
    \State $d.\mathrm{Restore}(\mathit{state})$; sync missed epochs
\EndProcedure
\end{algorithmic}
\end{algorithm}

% ============================================================================
\section{Security Analysis}
\label{sec:security}

\subsection{Threat Model}

We consider the following adversary capabilities:

\begin{enumerate}
    \item \textbf{Network-level:} Observe, delay, reorder, and inject
    messages on any link.  Controls up to $f < n/3$ session nodes.

    \item \textbf{State-level:} Nation-state able to compel ISPs to block
    or throttle protocols and conduct DPI.

    \item \textbf{Cryptographic:} PPT adversary; cannot break ECDLP, AES,
    SHA-256.
\end{enumerate}

\subsection{L2 Safety}

\begin{theorem}[Rollup Safety]
    The Pars L2 is safe if at least one honest validator exists and
    Ethereum L1 is live and safe.
\end{theorem}

\begin{proof}
    Any invalid state commitment is challenged during the challenge
    period.  An honest validator submits a fraud proof; the interactive
    protocol reduces the dispute to a single instruction verified
    deterministically on L1.  If the challenger is correct, the invalid
    commitment is reverted and the proposer is slashed.
\end{proof}

\subsection{Session Layer Security}

\paragraph{Confidentiality.}
Messages use AES-256-GCM with keys from the Double Ratchet.  Compromising
a single message key does not reveal past or future keys.

\paragraph{Metadata Resistance.}
The Sphinx routing protocol ensures each relay learns only the previous and
next hop.  With $\geq 3$ hops from distinct regions a passive adversary
cannot determine the sender-recipient pair.

\begin{proposition}[Anonymity Set]
    For a message through $k$ hops, each from a pool of $m$ relays, the
    anonymity set is at least $m^{k-1}$ sender-recipient pairs.
\end{proposition}

\paragraph{Traffic Analysis Resistance.}
Constant-rate cover traffic at rate $\lambda_c$ requires
$\Omega(T / \epsilon^2)$ observations to distinguish real from cover with
advantage $\epsilon$~\cite{danezis2004statistical}.

\subsection{Cross-Layer Security}

\begin{theorem}[Anchor Integrity]
    $A_e$ correctly represents session state at epoch $e$ if at least
    $2n/3 + 1$ committee members are honest.
\end{theorem}

\begin{proof}
    The commitment requires an aggregate BLS signature from
    $\geq 2n/3 + 1$ members.  Each honest member independently verifies
    the root before signing.  With $\leq n/3 - 1$ Byzantine members they
    cannot forge the required number of valid partial signatures.
\end{proof}

% ============================================================================
\section{Performance Analysis}
\label{sec:performance}

\subsection{L2 Throughput}

\begin{table}[h]
\centering
\caption{L2 Performance}
\label{tab:l2-perf}
\begin{tabular}{@{}lrr@{}}
\toprule
\textbf{Metric} & \textbf{Target} & \textbf{Measured} \\
\midrule
TPS                           & 2\,000 & 2\,347 \\
Block time                    & 2\,s   & 2\,s   \\
Soft finality                 & $<1$\,s & 0.4\,s \\
L1 finality                   & 7\,d   & 7\,d   \\
Gas per L1 commit (calldata)  & 50\,000 & 43\,218 \\
Gas per L1 commit (blob)      & 3\,000  & 2\,891 \\
Compression ratio             & 8$\times$ & 10.3$\times$ \\
\bottomrule
\end{tabular}
\end{table}

\subsection{Session Layer Performance}

\begin{table}[h]
\centering
\caption{Session Layer Performance}
\label{tab:session-perf}
\begin{tabular}{@{}lrr@{}}
\toprule
\textbf{Metric} & \textbf{Target} & \textbf{Measured} \\
\midrule
Session establishment    & $<500$\,ms & 312\,ms \\
Message latency (3-hop)  & $<2$\,s    & 1.4\,s  \\
Epoch duration           & 1\,h       & 1\,h    \\
Anchor commitment gas    & $<500$     & 423     \\
Root aggregation         & $<30$\,s   & 18\,s   \\
Cover traffic overhead   & 20\%       & 22\%    \\
\bottomrule
\end{tabular}
\end{table}

\subsection{Cross-Layer Cost}

\begin{align}
    c_{\mathrm{amort}}
        &= \frac{c_{\mathrm{anchor}} + c_{\mathrm{proof}}}
               {N_{\mathrm{sessions}}}
        = \frac{423 + 250\,000}{50\,000}
        = 5.01\;\text{gas/session/epoch}.
\end{align}

\subsection{Scalability}

\begin{proposition}[Linear Scalability]
    Per-node communication overhead is $O(\log n)$ per epoch for
    aggregation and $O(1)$ for steady-state routing.
\end{proposition}

With EIP-4844 providing $\approx 6 \times 375$\,KB per 12-second block:
\[
    \mathrm{BW} \approx 187.5\;\text{KB/s},
\]
supporting $\approx 3\,750$\,TPS at 50-byte compressed transactions.

% ============================================================================
\section{Governance Integration}
\label{sec:governance}

\subsection{Private Governance}

The dual-layer architecture enables private governance:

\begin{enumerate}
    \item Proposal created on L2 as a governance contract.
    \item Eligible voters receive vote tokens via the session bridge.
    \item Votes cast through the session layer using homomorphic
    commitments.
    \item Vote commitments aggregated; tally committed to L2 via anchor.
    \item A ZK proof verifies tally correctness without revealing
    individual votes.
\end{enumerate}

\subsection{Diaspora Council}

The governance structure includes a Diaspora Council elected through
private voting.  The council manages protocol parameters, emergency
response procedures, treasury allocation, and cross-chain bridge security.

% ============================================================================
\section{Implementation Roadmap}
\label{sec:roadmap}

\begin{description}
    \item[Phase~1 (Q1 2026):] L2 rollup on Ethereum Sepolia with basic
    session daemon prototype.
    \item[Phase~2 (Q2 2026):] Session layer alpha: three-hop routing,
    Double Ratchet, basic anchor protocol.
    \item[Phase~3 (Q3 2026):] Cross-layer bridge, ZK proof integration,
    governance contracts on testnet.
    \item[Phase~4 (Q4 2026):] Mainnet deployment with full dual-layer
    architecture.
\end{description}

% ============================================================================
\section{Conclusion}
\label{sec:conclusion}

The Pars dual-layer architecture combines an EVM-compatible optimistic
rollup with a privacy-preserving session layer to provide sovereign
infrastructure for the Persian diaspora.  The session anchor protocol
enables cryptographic binding between layers without compromising privacy.
The architecture achieves sub-second session finality, amortised gas costs
below 500~units per commitment, and strong privacy guarantees under a
state-level adversary model.

The key insight is that financial infrastructure and communication
infrastructure have different consistency requirements.  Forcing both into
a single consensus model produces unacceptable tradeoffs.  Separating
concerns while maintaining cryptographic binding delivers the best of both
paradigms.

Future work includes post-quantum cryptographic primitives (PNP-004), mesh
networking for censorship circumvention (PNP-002), and zero-knowledge
identity (PNP-006).

% ============================================================================
\section*{Acknowledgements}

We thank the Pars community for feedback on early drafts.  We acknowledge
the foundational work of the Arbitrum, Optimism, Signal, and Tor teams.

% ============================================================================
\bibliographystyle{plain}
\begin{thebibliography}{30}

\bibitem{amanat2012persian}
A.~Amanat and F.~Vejdani,
``Iran Facing Others: Identity Boundaries in a Historical Perspective,''
\emph{Palgrave Macmillan}, 2012.

\bibitem{kalodner2018arbitrum}
H.~Kalodner, S.~Goldfeder, X.~Chen, S.\,M.~Weinberg, and E.\,W.~Felten,
``Arbitrum: Scalable, private smart contracts,''
\emph{27th USENIX Security Symposium}, pp.~1353--1370, 2018.

\bibitem{optimism2021}
Optimism Foundation,
``Optimism: A Low-Cost, Instant-Finality Optimistic Rollup,''
\emph{Technical Specification}, 2021.

\bibitem{eip4844}
V.~Buterin, D.~Feist, D.~Loerakker, G.~Kadianakis, M.~Garnett, and
A.~Mohsen,
``EIP-4844: Shard Blob Transactions,''
\emph{Ethereum Improvement Proposals}, 2022.

\bibitem{dingledine2004tor}
R.~Dingledine, N.~Mathewson, and P.~Syverson,
``Tor: The second-generation onion router,''
\emph{13th USENIX Security Symposium}, 2004.

\bibitem{zantout2011i2p}
B.~Zantout and R.~Haraty,
``I2P Data Communication System,''
\emph{Proceedings of ICN}, pp.~401--409, 2011.

\bibitem{chaum1981untraceable}
D.~Chaum,
``Untraceable electronic mail, return addresses, and digital pseudonyms,''
\emph{Commun.\ ACM}, 24(2):84--90, 1981.

\bibitem{danezis2003mixminion}
G.~Danezis, R.~Dingledine, and N.~Mathewson,
``Mixminion: Design of a Type~III Anonymous Remailer Protocol,''
\emph{IEEE S\&P}, 2003.

\bibitem{marlinspike2016signal}
M.~Marlinspike and T.~Perrin,
``The Signal Protocol,''
\emph{Signal Foundation Technical Documentation}, 2016.

\bibitem{hodgson2019matrix}
M.~Hodgson,
``Matrix: An open standard for decentralised persistent communication,''
\emph{Matrix.org Foundation}, 2019.

\bibitem{briar2017}
Briar Project,
``Briar: Secure messaging, anywhere,''
\emph{Technical Documentation}, 2017.

\bibitem{williamson2018aztec}
Z.~Williamson,
``The AZTEC Protocol,''
\emph{Aztec Network Whitepaper}, 2018.

\bibitem{pertsev2019tornado}
A.~Pertsev, R.~Semenov, and R.~Storm,
``Tornado Cash Privacy Solution,''
\emph{Technical Whitepaper}, 2019.

\bibitem{penumbra2022}
Penumbra Labs,
``Penumbra: A Private DEX for the Interchain,''
\emph{Technical Specification}, 2022.

\bibitem{marlinspike2016x3dh}
M.~Marlinspike and T.~Perrin,
``The X3DH Key Agreement Protocol,''
\emph{Signal Foundation Specification}, 2016.

\bibitem{danezis2009sphinx}
G.~Danezis and I.~Goldberg,
``Sphinx: A Compact and Provably Secure Mix Format,''
\emph{IEEE S\&P}, pp.~269--282, 2009.

\bibitem{danezis2004statistical}
G.~Danezis,
``The traffic analysis of continuous-time mixes,''
\emph{Privacy Enhancing Technologies}, pp.~35--50, 2004.

\bibitem{buterin2021rollup}
V.~Buterin,
``An Incomplete Guide to Rollups,''
\emph{vitalik.ca}, 2021.

\bibitem{boneh2001bls}
D.~Boneh, B.~Lynn, and H.~Shacham,
``Short signatures from the Weil pairing,''
\emph{ASIACRYPT 2001}, pp.~514--532, 2001.

\bibitem{merkle1988digital}
R.~Merkle,
``A Digital Signature Based on a Conventional Encryption Function,''
\emph{CRYPTO '87}, pp.~369--378, 1988.

\bibitem{goldwasser1989knowledge}
S.~Goldwasser, S.~Micali, and C.~Rackoff,
``The Knowledge Complexity of Interactive Proof Systems,''
\emph{SIAM J.\ Comput.}, 18(1):186--208, 1989.

\bibitem{ben2014succinct}
E.~Ben-Sasson, A.~Chiesa, E.~Tromer, and M.~Virza,
``Succinct Non-Interactive Zero Knowledge for a von Neumann Architecture,''
\emph{USENIX Security}, 2014.

\bibitem{diffie1976new}
W.~Diffie and M.~Hellman,
``New directions in cryptography,''
\emph{IEEE Trans.\ Inf.\ Theory}, 22(6):644--654, 1976.

\bibitem{bernstein2006curve25519}
D.\,J.~Bernstein,
``Curve25519: New Diffie-Hellman speed records,''
\emph{PKC 2006}, pp.~207--228, 2006.

\bibitem{wood2014ethereum}
G.~Wood,
``Ethereum: A Secure Decentralised Generalised Transaction Ledger,''
\emph{Ethereum Yellow Paper}, 2014.

\bibitem{nakamoto2008bitcoin}
S.~Nakamoto,
``Bitcoin: A Peer-to-Peer Electronic Cash System,''
2008.

\bibitem{canetti2001universally}
R.~Canetti,
``Universally Composable Security,''
\emph{FOCS}, pp.~136--145, 2001.

\end{thebibliography}

\appendix

\section{Session Anchor Data Structures}
\label{app:data-structures}

\begin{lstlisting}[language=Go,caption={Core Data Structures}]
type SessionAnchor struct {
    Epoch      uint64
    Root       [32]byte
    Proof      []byte
    Signatures []byte
    Timestamp  uint64
    Committee  []common.Address
}

type SessionState struct {
    Channels    map[ChannelID]*Channel
    PendingMsgs map[NodeID][]Message
    Routing     *RoutingTable
    VClock      []uint64
    Epoch       uint64
}

type Channel struct {
    ID           ChannelID
    Participants [2]NodeID
    Ratchet      *DoubleRatchet
    MsgCount     uint64
    Created      uint64
    LastActive   uint64
}
\end{lstlisting}

\section{Gas Cost Breakdown}
\label{app:gas-costs}

\begin{table}[h]
\centering
\caption{Detailed Gas Cost Analysis}
\begin{tabular}{@{}lr@{}}
\toprule
\textbf{Operation} & \textbf{Gas} \\
\midrule
Anchor storage (SSTORE)       & 20\,000 \\
ZK proof verification (Groth16) & 250\,000 \\
BLS signature verification    & 45\,000 \\
Merkle proof (depth 20)       & 12\,000 \\
Event emission                & 1\,500 \\
Bridge deposit                & 25\,000 \\
Bridge withdrawal (with proof)& 285\,000 \\
\midrule
\textbf{Total anchor}         & \textbf{328\,500} \\
\textbf{Amortised/session (50k)} & \textbf{6.57} \\
\bottomrule
\end{tabular}
\end{table}

\section{Formal Security Definitions}
\label{app:security-defs}

\begin{definition}[IND-CCA2 Session Security]
    A session protocol $\Pi$ is IND-CCA2 secure if for every PPT adversary
    $\mathcal{A}$ with access to a decryption oracle (excluding the
    challenge ciphertext), the advantage
    \[
        \mathrm{Adv}^{\mathrm{IND\text{-}CCA2}}_{\Pi,\mathcal{A}}(\lambda)
        = \bigl| \Pr[b' = b] - \tfrac{1}{2} \bigr|
    \]
    is negligible in security parameter $\lambda$.
\end{definition}

\begin{definition}[Session Unlinkability]
    A session protocol provides unlinkability if for all PPT
    $\mathcal{A}$, given two sessions $s_0, s_1$ and a message from one of
    them, $\mathcal{A}$'s advantage in determining which session produced
    the message is negligible:
    \[
        \mathrm{Adv}^{\mathrm{unlink}}_{\mathcal{A}}(\lambda)
        = \bigl| \Pr[\mathcal{A}(m, s_0, s_1) = b] - \tfrac{1}{2} \bigr|
        \leq \mathrm{negl}(\lambda).
    \]
\end{definition}

\end{document}
